\subsection{Сбор данных для симуляции}

Очень важным фактором качественных рендеров является понимание оптических свойств рендерируемых предметов. Поэтому в контексте реалистичного рендеринга важную роль занимают работы, занимающиеся сбором этих свойств. Среди них особенно интересна статья \cite{Jacques_2013}, обобщающая множество независимых исследований человеческих тканей в этой области.

Непосредственно собираются медицинские данные для рендеринга в основном из КТ (выдаёт единицы Хаундсфилда, HU) и МРТ. Их особенность --- скалярность. Поэтому перед тем, как загонять их в движок, требуется некоторая подготовка. Для этого вводится функция переноса, чаще всего $\mathrm{HU} \to (RGBA)$, иногда добавляется $\nabla\mathrm{HU}$ для выделения границ органов. Одним из примеров подхода с градиентом и еще большим числом параметров в функции переноса является работа \cite{kroes2012exposure}.

Иногда удается сопоставлять и более сложные оптические свойства. Хорошим примером является работа \cite{verleker2014optical}, в которой из КТ извлекают коэффициенты поглощения, рассеивания и анизотропию (описывающую распределение рассеивания света), по которым симулируется поведение света, чтобы узнать попадает ли он в нужные зоны головного мозга.